% --- Current Employment
\begin{rubric}{Ongoing Role}

  \entry*[2021--Present]%
  {%
    \textbf{\href{https://www.linkedin.com/in/dvbeckwitt/}{Graduate Research Assistant}}, University of Missouri%
    \par \textit{Focus:} X-ray diffraction instrumentation, GIWAXS simulation, thin-film growth, and AI-assisted scientific-computing workflows.%
    \par \vspace{0.05em}%
    \begin{itemize}
      \item \textbf{Lab operations:} Operate and maintain custom-built in-house x-ray diffraction and reflectivity systems for materials characterization.
      \item \textbf{Forward modeling:} Built a \textbf{Python}-based GIWAXS framework to extract site occupancies, anisotropic Debye-Waller factors, mosaicity, and geometric parameters
      (\altlink{https://www.researchgate.net/publication/377264935_Quantifying_Orientational_Order_of_PbI2_van_der_Waals_Films_with_X-ray_Diffraction_using_an_Area_Detector}{APS Mar 2023}).
      \item \textbf{Defect analysis:} Extended GIWAXS methods to \textbf{model diffuse scattering} from stacking faults and \textbf{quantify defect densities}
      (\altlink{https://meetings.aps.org/Meeting/MAR24/Session/S64.3}{APS Mar 2024}).
      \item \textbf{Film growth:} Grew phase-controlled PbI$_2$ films via \textbf{Chemical Vapor Deposition} while validating polytype fractions
      (\altlink{https://doi.org/10.1021/acsami.3c14559}{ACS AMI 2023}).
      \item \textbf{Scientific ML and AI workflows:} Develop \textbf{PyTorch} CNN workflows trained on simulated GIWAXS detector images for structure-analysis screening. The use LLM-assisted agents and reusable AI skills for code scaffolding, parameter exploration, documentation, reproducibility checks, and validation planning against scattering constraints.
    \end{itemize}
  }

\end{rubric}
